\subsection{Traitement audiovisuel}

 dans la première sous partie : considérations technique : intégrer les contraintes du corpus et voir comment ils s’y inscrivent 
 
"However, while shotboundary detection is certainly useful, reliably being able to de-tect scene boundaries would be an incredible technical feat, whichwould create tremendous potential for semantic and narrative anal-ysis of video and film material" \footcite{bhargav2019}

Les tâches automatiques pouvant être appliquées à l’audiovisuel sont multiples avec d’un côté les tâches pouvant être appliquées à l’image, et de l’autre, les tâches pouvant être appliquées au son. Citons par exemple la reconnaissance de locuteurs, la segmentation, transcription et classification audio, la reconnaissance de visages, l’identification et la classification d'objets, la transcription de textes dans l’image...  

Nous nous concentrerons dans cette partie sur la mise en place d’un processus de segmentation temporelle de l’image en différentes séquences, puis de leur indexation. En effet, chez Skinsoft, le besoin de segmentation est issu d’un corpus d’archives audiovisuelles muettes, de rushes, que la segmentation permettrait de diviser en séquences thématiques afin de les rendre intelligibles. Bien que la segmentation est souvent la condition préalable à un objectif de segmentation, par exemple, la vidéo doit être segmentée avant de pouvoir en extraire le contenu automatiquement (personnes, objets...), il s’agit ici d’une fonctionnalité à part entière à laquelle pourra être combinée une indexation.  

Très important : la définition des séquences comme quoi ce sont des segments pas seulement syntaxiques mais surtout sémantiques : donc il faut non seulement les segmenter mais aussi en extraire le sens. : doit être le fil conducteur de ma partie  justifier davantage le choix des outils : Les institutions patrimoniales n'ont pas de GPU dédiés ni d'infrastructure cloud. Cela conditionne directement le choix des outils : un modèle léger comme PySceneDetect ou 2SDS sera préférable à ViViT ou TDSNet dans ce contexte, même si ces derniers sont plus performants. Attard et Seychell (2025) montrent que les architectures séparées et légères donnent de meilleurs résultats pratiques que les modèles multimodaux complexes 

les archives ne peuvent pas être envoyées vers des services cloud pour des raisons de droits de propriété intellectuelle et de sécurité des données. L'algorithme doit être déployé localement dans l'infrastructure de l'application Skinsoft. C'est une contrainte technique majeure qui doit apparaître dans ta sous-partie car elle élimine d'emblée certaines solutions (APIs cloud comme AWS Rekognition ou Google Video Intelligence) et oriente vers des solutions déployables en local. 

Nos recherches ont montré que les études sur le sujet restent cantonnées aux expérimentations scientifiques et qu'elles n'ont pas été appliquées, à notre connaissance, au domaine des humanités numériques. Nous établirons ainsi un état des lieux de ces expérimentations, afin de proposer la mise en place d'un tel processus dans un contexte de gestion de collections audiovisuelles patrimoniales.  

De plus, la segmentation des archives audiovisuelles n’est jamais mise en place de manière isolée, puisqu’elle est plutôt la première étape d’un processus global. Ainsi, nous intègrerons également un processus automatique afin de pouvoir les décrire et les indexer au sein d’un système de gestion des collections.  

Les chercheurs qui analysent les objets visuels s'attachent de plus en plus à l'analyse vidéo car la détection du mouvement aide d'autres tâches d'analyse automatique de détection et de reconnaissance d'objets.  \footcite{brostow2009}


\subsubsection{Machine learning et deep learning}

Les technologies appliquées à l'analyse de contenus vidéo et audio, comprenant la computer vision et la reconnaissance de discours automatique (ASR) sont désignées sous le terme de Audiovisual Processing (AVP). \footcite{noord2021}
aussi désigné sous les termes suivants : Audiovisual Processing technologies, Automatic Metadata Extraction, Metadata Generation Mechanisms, Machine Learning Applications for Visual/Audio Data Annotation.

L’avantage de ces technologies est le suivant : appliquées à un grand nombre de données, elles permettent de passer du "close" au "distant" viewing, c’est-à-dire d’une vue des données rapprochée, réduite, à une vue d’ensemble. \footcite{noord2021} Ainsi, l’automatisation permettrait de passer de l’analyse manuelle de certaines archives audiovisuelles, à une vue d'ensemble sur les données afin d'en observer les tendances globales. L'application du traitement audiovisuel permettrait donc de changer notre rapport aux archives audiovisuelles.  

L’application des technologies d’intelligence artificielle au domaine des images (images fixe ou vidéo) est désignée sous le terme de computer vision, ou de vision par ordinateur. Le terme définit un sous-domaine d’application de l’intelligence artificielle qui permet d’entraîner un système d’intelligence artificielle à l’analyse des images fixes ou en mouvement \footcite{stryker2025} La vision par ordinateur s’appuie largement sur des modèles de machine learning, l’apprentissage machine, qui désigne une branche de l'intelligence artificielle regroupant les techniques, méthodes et processus nécessaires afin de rendre une machine ou un système informatique plus intelligent. \footcite{zotero-item-584} Le machine learning permet d’entraîner des algorithmes capables d’apprendre à partir de données d’entraînement, pour ensuite prédire ou prendre des décisions sur de nouvelles données. \footcite{zotero-item-702} La vision par ordinateur s’appuie sur l’apprentissage machine précisément pour sa capacité à entraîner des systèmes à prédire. En effet, les corpus d’images fixes et d’images en mouvement sont des phénomènes complexes à appréhender pour des algorithmes, contrairement à des données statistiques et mathématiques. La nature de ces corpus les rend imprévisibles autant dans leur forme que dans leur contenu. Les algorithmes d’intelligence artificielle ne peuvent donc pas seulement réaliser une analyse prédictive de ce type de corpus, définie par des critères mathématiques. \footcite{mahony2020} Les corpus audiovisuels nécessitent une analyse prédictive, que les algorithmes sont entraînés à réaliser via des méthodes de machine learning. \footcite{mahony2020} 

De plus, la plupart de ces corpus sont d’une taille conséquente, forçant ainsi la mise en place d'algorithmes de prédiction, puisque l’entraînement machine ne peut pas être mis en place sur l’ensemble des données mais seulement sur un échantillon donné. 

Pour des tâches telles que la segmentation en différentes séquences, qui implique l’analyse du contenu de l’image, sont utilisés des modèles d’intelligence artificielle reposant sur des réseaux de neurones, et notamment des réseaux de neurones profonds. \footcite{norindr2023} La méthode de mise en place de ces réseaux de neurones profonds en plusieurs couches, est désignée sous le terme de deep learning, apprentissage profond, une sous-catégorie du machine learning. Le deep learning est basé sur des réseaux de neurones artificiels inspirés du fonctionnement du cerveau humain. Les neurones sont en réalité des cellules de calcul organisées en différentes couches qui interagissent les unes avec les autres afin de s’adapter aux données reçues, et de prendre une décision quelconque. Davantage entrainé que programmé, un modèle fondé sur l’apprentissage profond permet d’agir sans supervision humaine, améliorant ainsi le rapport cout-efficacité. \footcite{mahony2020}  Il suffit presque de fournir au modèle quelques images annotées, et le réseau de neurones permet au modèle de repérer lui-même les caractéristiques qu’il doit analyser sur de nouvelles données. \footcite{mahony2020} Par rapport aux techniques traditionnelles de vision par ordinateur, le deep learning permet d’obtenir une plus grande précision dans les tâches d'analyse audiovisuelle telles que la classification d'images, la segmentation, la détection d'objets. \footcite{mahony2020} 

Au sein des réseaux de neurones, le type de réseau spécifiquement employé dans le domaine de la computer vision pour la classification d’images et la reconnaissance d’objets, est celui le réseau de neurones convolutifs (CNN). Appliqués à des corpus d’images, les CNN peuvent par exemple analyser une image selon une hiérarchie de caractéristiques afin d’en extraire les couleurs et les contours. \footcite{brachmann2017}. Les taches que permettent les CNN sont plutôt bien adaptées à l’image fixe, puisqu’ils extraient des éléments fixes de l’image.  

Néanmoins, d’autres modèles issus du deep learning tels que les transformers s’avèrent davantage efficace sur des tâches dédiées à l’image en mouvement, dont fait partie la segmentation. Contrairement aux CNN, les transformers fonctionnent sur un mécanisme d’auto-attention qui leur permet de détecter les relations et dépendances entre les données. \footcite{bergmann2025} Tandis que les CNN analysent les données de manière sérielle, sans prendre en compte les interdépendances qui les régissent, les transformers examinent simultanément les données dans leur ensemble et décident de se concentrer sur certaines étapes de chacune d'entre elles. \footcite{bergmann2025} Les CNN ne permettent pas, par exemple, de discerner les “dépendances à longue portée” telles que les corrélations entre les pixels entre différentes images qui ne seraient pas voisins. \footcite{bergmann2025} Ainsi, ils seraient davantage adaptés à l’analyse d’images en mouvement puisqu’ils prendraient en compte les relations entre les images, tandis que les CNN seraient davantage adaptés aux images fixes.  Idéalement, la mise en place d’une segmentation automatique des archives audiovisuelles doit donc être fondée sur une architecture de type transformers.  

Toutefois, dans un contexte patrimonial, avec des ressources limitées, ces algorithmes fondés sur l’intelligence artificielle possèdent leur lot de limites. Premièrement, les réseaux de neurones sont conçus comme des boites noires, ne permettant pas d’interpréter les analyses qui sont produites. \footcite{brachmann2017}  En effet, une fois entraînés, il est difficile d’en modifier les paramètres puisqu’ils sont entremêlés entre les différentes cellules de calcul. La réflexion produite est ainsi difficile à comprendre et à exposer. \footcite{mahony2020} Deuxièmement, les algorithmes de type transformers sont complexes et restent difficiles à adapter à un corpus audiovisuel lui-même complexe. \footcite{yuan2025}. Ils demandent également de grandes puissances de calcul et un temps de paramétrage important, ce qui les rend inadaptés à des contextes de gestion de collections patrimoniales. De plus, ils requièrent un grand nombre de paramètres issu d’un large jeu de données d’entraînement afin d’éviter un trop grand nombre de biais.  \footcite{pattichis2023} 
%À inclure ? sur les limites car modèles entrainés sur des corpus inadaptés : prendrait + de temps en validation et correction, pas stable
les modèles de détection neuronaux sont entrainés sur des datasets de clips courts. Or, la réalité auiovisuelle est bien différente. \footcite{pattichis2023} Des approches algorithmiques non neuronales ou hybrides seraient donc plus adaptées car fondée sur des calculs et des seuils matéhmatiques : plus stable. 





% TODO: \usepackage{graphicx} required
\begin{figure}
	\centering
	\includegraphics[width=0.7\linewidth]{images/screenshot001}
	\caption{mahony2020}
	\label{fig:screenshot001}
\end{figure}




Les approches hybrides : un algorithme de vision par ordinateur peut détecter des caractéristiques qui peuvent ensuite être transmies à des réseaux de neurones profonds pour la classification. 

Mais tout dépend de la tâche à effectuer : peut être effectué à partir de la computer vision si c'est juste un seuil à calculer.


transition vers la prochaine sous partie : parler des données d'entrainement (cf fichier IAvideo pattichis)

Toutes les tâches automatiques qui peuvent être appliquées à la vidéo nécessitent un entraînement multimodal : 
entrainement multimodal :  When we refer to multimodal or cross-modal learning, we are referring to learning using video, audio and/or text. These approaches generalize to a greater variety of visual tasks including action step localization, action segmentation, text-to-video retrieval, video captioning. \footcite{schiappa2023}


\subsubsection{Outils et modèles disponibles}

Différents projets ont tenté de mettre en place des algorithmes de segmentation vidéo. Bien que ces projets ne soient pas appliqués à des archives audiovisuelles, mais des corpus de vidéos moins spécifiques, il est tout de même intéressant d’en analyser la structure pour éventuellement les implémenter dans un contexte patrimonial. De plus, il est intéressant de comprendre sur quelles technologies s’appuient ces algorithmes. En effet, face aux limites de sréseaux euronaux et des algorithmes fondés sur le mécanisme d'auto attention, nous tenterons d'analyser des architectures plus simples et également plus fiables afin de les appliquer aux archives. 
Les modèles de détection neuronaux sont entrainés sur des datasets de clips courts. Or, la réalité audiovisuelle est bien différente. \footcite{pattichis2023} Des approches algorithmiques non neuronales ou hybrides seraient donc plus adaptées car fondée sur des calculs et des seuils matéhmatiques : plus stable. 

Nous analyserons la structure de ces algorithmes sans en mentionner les détails mathématiques et techniques puisque notre objectif est d’en proposer une implémentation métier, au sein d’un système de gestion des collections. Nous analyserons davantage leur adaptabilité aux données et leur capacité d’entrainement sur des donnés complexes.  

Une des solutions, afin de pouvoir implémenter des algorithmes plus légers, serait de ramener l'image ne mouvement à l'image fixe. En effet, le cinéma est généralement caractérisé par l'enchaînement de 24 images par seconde afin de produire une illusion du mouvement pour l'oeil humain. Nous pouvons éventuellement imaginer un outil qui puisse découper le mouvement en des images fixes qui le composent afin de pouvoir les analyser plus facilement, au lieu d'implémenter une architecture complexe fondée sur la compréhension des relations et du mouvement entre les images. Les images fixes seraient ensuite regroupées en séquences après uen analyse de leurs différence. Grâce à ce procédé, nous pourrions finalement appliquer les mêmes procédés d'analyse que sur l'image fixe, relativement plus légers, fondés sur la vision par ordinateur. 

C’est ce que fait la bibliothèque python pyscenedetect  avec différentes méthodes de détection. \footcite{zotero-item-706} Après, les plans sont regroupés de manière thématique après une analyse entre leurs différences : plus simples. 

Dans leur article, Xin présente les principes de segmentation appliqué à la vidéo, pronant le découpage syntaxique préalable à l'analyse sématique par des modèles d'intelligence artificielle. \footcite{xin2023}. Ainsi, il faudrait détecter les changements entre les scènes d'une vidéo en analysant préalablement les différences entre les images fixes qui la composent. Si deux images sont jugées suffisamment différentes selon un seuil prédéfini, elles appartiennent à des scènes distinctes. \footcite{xin2023}. Une ou plusieurs images clés pour chaque séquence sont retenues, puisque l'algorithme élimine les images en double, proches, et retient celle qui est représnetative de la présence des éléments dans l'image (data selection). \footcite{xin2023} Ce principe permet d'alléger nettement la puissance de calcul et la structure des algorithme. Ces imges clés seraient ensuite passées dans un algorithme fondé sur un réseau neuronal pour effectuer une analyse sémantique. 

Pyscene Detect est un outil opensource et en libre accès fondé sur des tâches de computer vision classique, qui matérialise concrètement ce principe de comparaison première des images fixes. Il utilise des détecteurs mathématiques pour analyser les différences entre les images tirées d'une vidéo. \footcite{zotero-item-708} Il s’agit ainsi d’un outil basé sur des seuils et mesures visuelles mathématiques  et paramétrables qui ne nécessitent pas d’entraînement particulier. L'outil permet donc d'isoler les images fixes de l'image animée, de les regrouper par similarité et de les dissocier selon leurs différences. Trois méthodes de détection sont proposées par l’outil : une détection des fondus entrants et sortants qui repère les niveaux de noir dans l’image, une détection de coupes rapides qui repère les changements de contenu entre images, une détection de coupes fondée sur la différence d’un plan avec les plans adjacents. \footcite{zotero-item-708} La méthode utilisée dépend de la nature du corpus analysée. Si les changements entre les séquences sont peu marqués, peu évidents, la méthode detect-adaptative sera plutôt employée. Une fois la méthode choisie, il est possible de générer un fichier statistique afin de déterminer les paramètres à appliquer. Il est notamment possible de paramétrer la valeur du seuil voulu, selon la détection de changement souhaitée entre les différents plans. \footcite{zotero-item-708} La méthode de détection fondée sur le contenu analyse la différence entre chaque paire d’images adjacentes et déclenche un changement de scène lorsque cette différence dépasse la valeur seuil. Cette valeur doit être très faible entre des images similaires et augmenter lorsqu'un changement important de contenu est détecté. La valeur seuil doit être définie de manière à ce que la plupart des scènes se situent en dessous de cette valeur, tandis que les scènes où des changements se produisent doivent la dépasser (déclenchant ainsi un changement de scène). \footcite{zotero-item-708} L'outil retourne ensuite des marqueurs temporels associés à la vidéo correspondant aux coupures détectées. 
Un détail important de l'outil est sa capacité à enregistrer automatiquement une image clé pour chaque scène détectée, associée au fichier vidéo d'origine. C'est cette iage clé qui sera réemployée au sein du processus suivant, celui de l'analyse sémentique de la séquence, afin de ne pas transmettre une multitude d'image auux algorithmes. 
Toutefois, certaines mises en application de l'outil soulignent ses limites. Bhargav expique par exemple que l'outil détecte beaucoup plus de coupures de plans que nécessaire au sein d'archives du cinéma des premiers temps. \footcite{bhargav2019} Un problème qui peut être lié à la qualité amoindrie de ces archives, même numérisées.  \footcite{bhargav2019}. En effet, selon Norindr, la qualité détermine la performance de l'outil utilisé et des résultats reçus sur l'analyse de l'image. \footcite{norindr2023} De plus, Bhargav souligne la nécessité de savoir comment paramétrer l'outil de manière adaptée au corpus analysé. \footcite{bhargav2019} Toutefois, un des avantages majeurs de l'outil pyscene detect est sa disponibilité et la possibilité de le mettre en oeuvre facilement et gratuitement. De plus, c'est un outil fondé sur des tâches classiques de vision par ordinateur analytiques qui n'ont pas besoin d'entraînement particulier sur des données annotées. Son paramétrage est également documenté et facilement modifiable, contrairement à ceux des réseaux de neurones. 



Après pyscenedetect : apposer une analyse sur les séquences. Pas adapté pour des archives plus complexes. Sur notre corpus : ancien donc ruptures moins visibles avec la dégradation. Pourrait convenir, mais pour d’autres corpus, appliquer un pipleine basé sur des mécanismes plus lourds d’auto-attention.

Quand déclencher le processus ??? réalité métier : application de l’IA : à la demande du client ? Sur certaines archives ? À l’ingest d’un fichier numérique ? Sur tous les fichiers numériques ? 

\subsection{Mise en application}

Mettons à l'épreuve l'adaptabilité de l'outil pyscenedetect pour la segmentation des vidéos en séquences, au corpus d'archives audiovisuelles géré au sein du logiciel de Skinsoft. Ce corpus, rappelons-le, est composé d'archives issues des années 1910 à 1920 essentiellement documentaures, à plans fixes, en noir et blanc et muettes. Nous imaginerons ici ce que pourrait produire l'outil sur un tel corpus et quelle seraient les données générées par le processus, avant même de réfrléchir à une interface entre l'outil et l'utilisateur. Le corpus pris comme exemple se caractérise par son ancienneté, et donc son indéniable dégradation de la qualité de l'image. L'outil pourrait rencontrer des difficulté afin d'identifier les changements de plans et les coupures. Toutefois, la structure de ces archves en une succession de plans fixes, dont les coupures sont marquées, en fait un corpus ralitivement bien adapté à l'outil, qui pourrait identifier les coupes plus facilement. 
Par ailleurs, pyscenedetect est un outil davantage utilisé pour la détection de plans différents au sein d'une image animée. Or, dans ce mémoire, nous cherchons à savoir comment segmenter une archives audiovisuelle en séquences. Dès lors, comment paramétrer un tel outil pour la détection d'éléments syntaxiques plus larges que ceux des plans ? 
Une des solutions est d'augmenter la valeur du seuil de coupures (treshold), pour que l'outil ne détecte plus des micro coupures entre les images mais des coupes plus larges, à l'échelle de séquences. \footcite{zotero-item-708} L'outil enregistrerait un changement seulement lors d'un changement radical d'espace temps à l'image, losrque l'image change significativement. Ce paramètre reste limité puisque plusieurs plans peuvent être radicalement différents mais appartenir à la même séquence. Par exemple, un gros plan et  un plan d'ensemble. Dans le cas du corpus d'archives décrit, les plans restent larges pour la grande majorité, montrant l'ensemble du décor dans lequel ils s'inscrivent. 
Par ailleurs, il est possible de paramétrer une durée minimale de scène à détecter avec le paramètre "min-scene-len", pour diriger la coupure seulement après un laps de temps défini. \footcite{zotero-item-708}. Ce paramètre suppose une analyse statistiue du corpus afin de définir la durée moyenne d'une séquence au sein des archives. L'institution conservatrice de ce corpus documente justement déjà les séquences sur ses archives audiovisuelles à l'aide de time codes pour en déterminer les séquences. Une analyse statistique comparative pourrait être réalisée sur ces données documentées de time codes afin de déterminer la durée moyenne d'une séquence au sein du corpus. 
Enfin, sur ce corpus, le mode de détection adaptatif pourrait être privilégié, couplé à une analyse d'images plus importante. \footcite{zotero-item-708} Ainsi, il est possible de paramétrer le mode de détection de telle manière à ce que l'outil compare le contenu visuel d'une image par rapport à une plus grande variété des images qui l'entourent. Ce paramétrage permettrait à l'outil de détecter uniquement les changements majeurs dans l'image, et donc les séquences. 
Une fois ces paramètres définis, le processus une fois enclenché produit deux types de données. Il produit un fichier json ou csv contenant les time codes de début et de fin identifiés comme séquences au sein de chaque fichier vidéo ingesté. Chaque séquence possède donc un identifiant unique associé à des times codes de début et de fin. L'outil permet également d'extraire pour chaque séquence un fichier image, qui correspond à l'image fixe médiane, censée être la plus représentative de la séquence et que nous désignerons sous le terme d'image clé. \footcite{zotero-item-708} 

Toutefois, si ce paramétrage permettrait une plus grande précision dans la détection de séquences, il comporte des limites. Ces limites sont dûes à la définition d'une séquence, qui ne se définit pas seulement comme ue caractéristique syntaxique au cinéma, comme c'est le cas du plan, mais bien par une double caractéristique syntaxique et sémantique. En effet, un changement de séquence est un choix narratif, un choix de sens au sein du flux animé. C'est pourquoi un outil fondé sur des seuils mathématiques appliqués au contenu visuel comporte des limites. Nous proposons alors de réaliser une analyse sémantique sur les séquences détectées, afin de ré associer la syntaxe au ens, qui en fait une séquence, en réutiliant les données produites par le processus : les times codes et l'image clé.


Il est courant d’utiliser des modèles préentraînés sur un grand jeu de données et de le réajuster ensuite au jeu de données en question, au sein d’un contexte précis. \footcite{norindr2023} Ce modèle peut justement servir de modèle de caractéristiques à analyser sur les vidéos, un second entraînement, dans le contexte peut aider à l’affiner et à l’adapter au corpus spécifique. Le rendre d’abord familier au langage audiovisuel. 








------------------------------ 

Malgré de grandes avancées dans le processing audiovisuel, avec des outils comme la computer vision appliquée à la vidéo, et la reconnaissance de discours automatique (ASR), ces outils restent largement inaccessibles et hors de portée pour les institutions patrimoniales. \footcite{noord2021}. Ils sont utilisés par des entreprises commerciales pour du montage vidéo par exemple. Le manque d'utilisation de ces outils dans le domaine des sciences humaines expliquent l'inexistance de données d'entraînement adaptées au secteur. \footcite{noord2021} 
l'analyse vidéo par un processus humain prend un temps considérable, puiqu'il faut lire le fichier audiovisuel et en analyser les composants, le langage. Ainsi, l'indexation des archives audiovisuelles se cantonne souvent aux métadonnées, sans en exposer les données, le contenu.

Et donc en pratique, comment appliquer le traitement audiovisuel à des corpus d'archives audiovisuelles patrimoniales ? 






\subsubsection{entraîner ces outils et ces modèles}

les modèles de détection neuronaux sont entrainés sur des datasets de clips courts. Or, la réalité auiovisuelle est bien différente. \footcite{pattichis2023} Des approches algorithmiques non neuronales ou hybrides seraient donc plus adaptées car fondée sur des calculs et des seuils matéhmatiques : plus stable. 

détection des coupures : 
approche algorithmique non neuronale (2SDS, pyscenedetect), (xin et attard) pipeline technique en 4 étapes (down sampling (miniaturisation), conversion en niveaux de gris, calcul d'empreinte numérique par hachage binaire des pixels, et calcul de la distance de Hamming pour marquer la rupture visuelle (xin2023)

efficace sur les plans fixes et mvts lents 

mettre en rapport avec les travaux d'Arnold  Tilton (2019), qui utilisent une méthode similaire combinant différences d'histogrammes et réduction de résolution des pixels (downsampling) pour mesurer les ruptures de palettes de couleurs et de luminosité d'une frame à l'autre, complétée par des règles logiques pour éviter les faux positifs lors des fondus (fade-in/fade-out).

L'apport de l'IA (Vers la classification sémantique) : La détection brute des coupures ne suffit plus ; l'enjeu actuel est la Shot semantics (décrire le contenu d'un plan et sa typologie) (Arnold  Tilton, 2019).

Le rôle des modèles (CNN et Détecteurs) : Pour analyser le contenu de ces plans, on utilise l'apprentissage supervisé (Supervised Learning), où des algorithmes apprennent à partir de données labellisées par l'humain (DeLua, 2024). C'est le cas des réseaux de neurones convolutifs (CNN) comme Faster R-CNN ou YOLO (un détecteur de type Single Shot qui prédit et classifie "en une seule fois") pour la détection d'objets ou de visages (Arnold  Tilton, 2019 ; Robert, 2026).

Le méta-algorithme de classification : Arnold  Tilton (2019) montrent qu'en combinant la position des visages détectés et les coupures de plans, on peut créer un méta-algorithme capable de classifier automatiquement le type de plan (gros plan/close-up, plan amorce/over the shoulder), permettant de comprendre comment le cadrage construit le récit.

\subsubsection{analyse du mouvement}

au lieu d'utiliser des larges modèles qui tentent de tout classifier, plutot utiliser des modèles "low parameter", c'est-à-dire des modèles 3D-CNN légers, avec moins de paramètres dédiés à la reconnaissance d'une seule activité précise. prenennt moins de mémoire vidéo. \footcite{pattichis2023}

stratégie de réduction de l'annotation humaine : slf supervised learning, semi supervised, active learning, zero shot \footcite{pattichis2023}

confronter les architectures de réseaux de neurones (CNN vs. Transformers), introduire l'apprentissage non supervisé (Unsupervised Learning / DeepCluster),

Les limites des CNN 2D classiques (comme ResNet) : ils analysent la vidéo comme des images discrètes et indépendantes, ignorant la continuité du mouvement \footcite{xin2023}. Se tourner vers les vision transformers : CLIP, ViViT (bergmann et al), mécanismes d'auto attention qui analysent l'intégralité d'une séquence

Les modèles de Video Semantic Segmentation (VSS) et le dilemme entre les méthodes orientées "vitesse" (approche par keyframes peu précise) et orientées "précision" (TDNet, flux optique, mécanismes d'auto-attention) \footcite{zhou2022}.

La gestion du coût de calcul : présentation de TDSNet (TFRM pour extraire les zones dynamiques par soustraction d'images consécutives, et MMRM pour pondérer selon la magnitude réelle du mouvement) \footcite{yuan2025}. règle ce dilemme 

alternatives non supervisées : la supervision demande des données annotées (caron et al), l'apprentissage non supervisé (Unsupervised Learning) comme DeepCluster (Caron et al., 2018). Ce type de modèle regroupe les images ou frames similaires en "clusters" algorithmiques (via K-means) pour s'auto-entraîner sans intervention humaine (Caron et al., 2018). Si DeepCluster s'avère robuste face à des changements de distribution d'images (ce qui est prometteur pour des archives spécifiques), il présente deux limites majeures : il introduit des biais imprévisibles en remplaçant les labels par des métadonnées brutes (inacceptable pour la rigueur historienne) et a été conçu sur des bases gigantesques (ImageNet, 1.28 million d'images), le rendant difficilement transposable à des collections d'archives réduites (Caron et al., 2018 ; Norindr, 2023).

infrastructure locale : amener l'algorithme aux données : C'est la discussion centrale de votre partie. Les outils d'AVP commerciaux (souvent basés sur le Cloud) sont inutilisables pour les archives en raison des droits de propriété intellectuelle et de la sécurité des données (van Noord et al., 2021). La solution réside dans des architectures locales comme DANE (projet CLARIAH), basées sur un système décentralisé de workers (van Noord et al., 2021). Ce type d'infrastructure produit des métadonnées exploitables, mais aussi des vecteurs de caractéristiques (représentations abstraites) ouvrant la voie à la recherche par similarité visuelle ou sonore (van Noord et al., 2021).

La note finale sur l'explosion des données : L'indexation automatique par IA (comme avec DigInPix ou les modèles de type VLM) génère une masse d'informations locales phénoménale : les bases de données constituées peuvent être plus lourdes en taille que les fichiers vidéos d'origine (Moiraghi  Letessier, 2018).

transition : nécessité d'un mapping vers le FRBR : faire correspondre la sortie de ces pipelines avec un modèle de données. 